\chapter{Simulation from Python}
\label{ch:simulation}

This chapter uses a small linear pathway throughout: a species \code{S1} made from a fixed source \code{X0},
converted to \code{S2}, and \code{S2} removed.

\section{Loading a model}

A model is SBML text. Antimony, a short human-readable format, is converted to SBML by
\code{antimony\_to\_sbml}:

\begin{lstlisting}[language=Python]
import roadrunner_nitro as rn

sbml = rn.antimony_to_sbml('''
    J1: $X0 -> S1; k1*X0     # S1 made from X0 ($: X0 is held fixed)
    J2: S1 -> S2;  k2*S1
    J3: S2 -> ;    k3*S2
    X0 = 10; k1 = 0.5; k2 = 0.2; k3 = 0.4
''')
r = rn.RoadRunner(sbml)            # or rn.RoadRunner(open('model.xml').read())
\end{lstlisting}

Loading compiles the model to native code, which takes milliseconds even for large models. The lists of the
model's elements are on \code{r.model}:

\begin{lstlisting}[language=Python]
r.model.getFloatingSpeciesIds()    # ['S1', 'S2']
r.model.getReactionIds()           # ['J1', 'J2', 'J3']
r.model.getGlobalParameterIds()    # ['k1', 'k2', 'k3']
\end{lstlisting}

\section{Time courses}

\code{simulate(start, end, points)} integrates the model and returns a named array, one row per time point. The
columns are the \emph{time-course selections}: by default time and the concentrations of the floating species.
Any species, parameter, reaction rate or rule variable can be selected:

\begin{lstlisting}[language=Python]
r.timeCourseSelections = ['time', 'S1', 'S2', 'J2']
result = r.simulate(0, 20, 5)
print(result)
\end{lstlisting}
\begin{lstlisting}[style=out]
    time,      S1,      S2,      J2
 [[    0,       0,       0,       0],
  [    5,  15.803, 4.99471,  3.1606],
  [   10, 21.6166, 9.34557, 4.32332],
  [   15, 23.7553, 11.2863, 4.75106],
  [   20, 24.5421, 12.0463, 4.90842]]
\end{lstlisting}

The model is left at the end time: a second \code{simulate} continues from there unless you reset it.
\code{simulate(times=[...])} reports exactly the times listed.

\subsection{Reading and changing values}

Values are read and set by id, as with a dictionary:

\begin{lstlisting}[language=Python]
r['S1']             # 24.542..., the current value
r['k2'] = 0.3       # change a parameter
r.reset()           # back to the initial state; parameter changes are kept
r.resetAll()        # initial state and the model's own parameter values
\end{lstlisting}

\code{getValue(id)} and \code{setValue(id, value)} do the same. For a species, \code{'S1'} is the
concentration and \code{'[S1]'} names it explicitly.

\subsection{Many time courses at once}

\code{simulateBatch} runs one time course per parameter set, in parallel on every core, and leaves the model
unchanged. It returns an array of shape (sets, points, selections):

\begin{lstlisting}[language=Python]
out = r.simulateBatch({'k1': [0.25, 0.5, 1.0]}, 0, 20, 5, selections=['time', 'S2'])
out.shape            # (3, 5, 2)
out[:, -1, 1]        # S2 at time 20 for each k1: [ 6.48, 12.5, 24.55]
\end{lstlisting}

\subsection{Integrators}

The default integrator (\code{'cvode'}) switches between a variable-order BDF method for stiff problems and a
Rosenbrock method, and locates events and discontinuities exactly. \code{r.setIntegrator(name)} selects another:
\code{'gillespie'} (stochastic simulation), \code{'rk4'}, \code{'rk45'} or \code{'euler'}. Settings use
libroadrunner's names:

\begin{lstlisting}[language=Python]
r.integrator.relative_tolerance = 1e-10
r.integrator.absolute_tolerance = 1e-14

r.setIntegrator('gillespie')       # stochastic runs, reproducible with a seed
r.setSeed(1)
r.simulate(0, 20, 5)
\end{lstlisting}

\section{Steady states}

\code{steadyState()} moves the model to its steady state. It uses Newton's method on the system reduced by its
conservation laws, and if Newton's method fails from the current state it integrates towards the steady state
and tries again, so it finds steady states that Newton's method alone misses.

\begin{lstlisting}[language=Python]
r.steadyState()
r.getSteadyStateValues()        # [25. , 12.5]  (S1, S2)
r.model.getReactionRates()      # [5., 5., 5.]
\end{lstlisting}

\section{Jacobians and control analysis}

Matrices come back as named arrays, with \code{rownames} and \code{colnames}:

\begin{lstlisting}[language=Python]
J = r.getFullJacobian()
print(J)
\end{lstlisting}
\begin{lstlisting}[style=out]
        S1,   S2
S1 [[ -0.2,    0],
S2  [  0.2, -0.4]]
\end{lstlisting}

Control coefficients and elasticities are computed at the steady state, with exact derivatives:

\begin{lstlisting}[language=Python]
r.getCC('J3', 'k1')     # 1.0: scaled flux control coefficient of k1 on J3
r.getuCC('S2', 'k2')    # 0.0: S2 = k1 X0 / k3 does not depend on k2
r.getEE('J2', 'S1')     # 1.0: scaled elasticity of J2 with respect to S1
r.getScaledFluxControlCoefficientMatrix()
\end{lstlisting}

Other matrices follow libroadrunner's names, for example:

\begin{lstlisting}[language=Python]
r.getFullStoichiometryMatrix()
r.getLinkMatrix()
r.getConservationMatrix()
r.getReducedJacobian()
r.getFullEigenValues()
r.getScaledElasticityMatrix()
r.getScaledConcentrationControlCoefficientMatrix()
\end{lstlisting}
